\section{The square root of two is irrational}\label{ch04:sec:sqrt2}
Recall from Section~\ref{ch01:sec:notation} that a \emph{rational number} can be written as a fraction $a/b$ of integers with $b\ne0$; a real number that is not rational is called \emph{irrational}. The symbol $\sqrt2$ denotes the nonnegative real number whose square is $2$.

\par\noindent\begin{minipage}{\linewidth}
\begin{theorem}\label{thm:sqrt2}
The number $\sqrt2$ is irrational.
\end{theorem}
\end{minipage}\par

Why argue by contradiction? The theorem says what $\sqrt2$ is \emph{not}, and a direct proof would have to rule out every fraction at once. Supposing the opposite hands us one fraction of integers equal to $\sqrt2$, with a numerator and denominator to calculate with. Statements that no object of some kind exists are often proved by contradiction for this reason.

The proof chooses a fraction for $\sqrt2$ with the smallest possible positive denominator. For an ordinary rational number such as $0.75=3/4=6/8=9/12=\cdots$, the possible positive denominators are $4,8,12,\ldots$, and choosing the smallest picks out $3/4$, which cannot be reduced further.

\noindent\textit{Plan.} Suppose $\sqrt2=a/b$ with $b$ as small as possible. Squaring will force both $a$ and $b$ to be even. Then cancelling a factor of two gives a fraction for $\sqrt2$ with a smaller positive denominator, which contradicts the choice of $b$.
\begin{proof}
Suppose, for contradiction, that $\sqrt2$ is rational: $\sqrt2=p/q$ for some integers $p$ and $q$ with $q\ne0$. If $q$ is negative, we may rewrite $p/q$ as $(-p)/(-q)$, whose denominator is positive. So the set of positive integers $j$ for which $\sqrt2=a/j$ with some integer $a$ is nonempty. By the well-ordering property (Section~\ref{ch04:sec:contrapositive}), it has a smallest member; call it $b$, and let $a$ be an integer with $\sqrt2=a/b$.

Squaring gives $2=a^2/b^2$, and multiplying by $b^2$ gives $a^2=2b^2$. Since $b^2$ is an integer, $a^2$ is even, so $a$ is even by Lemma~\ref{thm:parity-square}; write $a=2c$ with $c$ an integer. Substituting into $a^2=2b^2$ gives
\[
(2c)^2=2b^2,\qquad\text{that is,}\qquad 4c^2=2b^2,\qquad\text{so}\qquad b^2=2c^2.
\]
Thus $b^2$ is even, and Lemma~\ref{thm:parity-square} shows that $b$ is even; write $b=2d$ with $d$ an integer. Since $b$ is positive, so is $d=b/2$, and $d<b$ because $b=d+d$. Cancelling the common factor $2$ gives
\[
\sqrt2=\frac{a}{b}=\frac{2c}{2d}=\frac{c}{d}.
\]
So $\sqrt2=c/d$ with $c$ an integer and $d$ a positive integer smaller than $b$, contradicting the choice of $b$. Therefore $\sqrt2$ is irrational.
\end{proof}

Parity alone gives no contradiction: fractions with even numerator and denominator, such as $6/8$, certainly exist. The contradiction comes from parity together with the choice of $b$.

No computation could replace this proof. Computers have produced billions of digits of $\sqrt2\approx1.41421356$, but any finite list of digits is a fraction, such as $141421356/100000000$. Only a proof can show that \emph{no} fraction of integers, however large, equals $\sqrt2$.

\pauseq{Where did the proof use the fact that the denominator $b$ is positive?}{In two places. First, the well-ordering property applies only to sets of \emph{positive} integers, and negative denominators have no smallest one: from $\sqrt2=a/b$ with $b>0$ we also get $\sqrt2=(-a)/(-b)=(-2a)/(-2b)=\cdots$, with denominators $-b,-2b,\ldots$ decreasing without end. Second, we needed $d=b/2$ to be smaller than $b$, which fails for negative $b$: if $b=-4$, then $b/2=-2>b$.}
